% =============================================================================
% Roteiro: como adicionar um modulo ao Morphe -- o exemplo da soma de dois
% sinais. Arquivo unico, compila no Overleaf com pdfLaTeX.
%
% O codigo mostrado aqui e o da branch estagio/modulo-soma do repositorio
% (Quartus/soma.v, C/morphe_server.c, Python/soma_window.py ...).
% =============================================================================
\documentclass[11pt,a4paper,oneside]{article}

\usepackage[utf8]{inputenc}
\usepackage[T1]{fontenc}
\usepackage{lmodern}
\usepackage[brazilian]{babel}
\usepackage[a4paper,margin=2.5cm]{geometry}
\usepackage{microtype}
\setlength{\emergencystretch}{3em}
\usepackage{xcolor}
\usepackage{booktabs}
\usepackage{tabularx}
\usepackage{array}
\usepackage{enumitem}
\usepackage[most]{tcolorbox}
\usepackage{listings}
\usepackage{tikz}
\usetikzlibrary{positioning,arrows.meta,fit,backgrounds}
\usepackage[hidelinks]{hyperref}

% --- cores (as do manual: azul = cliente, verde = HPS, laranja = FPGA) ---
\definecolor{cliente}{HTML}{2563EB}
\definecolor{hps}{HTML}{15803D}
\definecolor{fpga}{HTML}{C2410C}
\definecolor{fundo}{HTML}{F5F7FA}
\definecolor{comentario}{HTML}{6B7280}

% Nome de arquivo ou comando, sempre em pé.
\newcommand{\arq}[1]{{\upshape\ttfamily\detokenize{#1}}}
\newcommand{\opt}[1]{{\upshape\ttfamily -{}-\detokenize{#1}}}

% --- caixas ---
\newtcolorbox{atencao}{colback=orange!6,colframe=orange!70!black,
  title=Atenção,fonttitle=\bfseries\small,fontupper=\small,boxrule=0.5pt,arc=2pt}
\newtcolorbox{pronto}{colback=green!5,colframe=green!45!black,
  title=Pronto para reaproveitar,fonttitle=\bfseries\small,fontupper=\small,
  boxrule=0.5pt,arc=2pt}
\newtcolorbox{confira}{colback=blue!4,colframe=blue!50!black,
  title=Como saber que deu certo,fonttitle=\bfseries\small,fontupper=\small,
  boxrule=0.5pt,arc=2pt}

% --- código ---
\lstset{
  basicstyle=\ttfamily\footnotesize,
  backgroundcolor=\color{fundo},
  commentstyle=\color{comentario},
  keywordstyle=\color{cliente}\bfseries,
  stringstyle=\color{fpga},
  frame=single, framerule=0pt, framesep=4pt,
  xleftmargin=6pt, xrightmargin=6pt,
  breaklines=true, breakatwhitespace=true,
  columns=fullflexible, keepspaces=true,
  showstringspaces=false, tabsize=4,
  upquote=true,
  literate={á}{{\'a}}1 {é}{{\'e}}1 {í}{{\'i}}1 {ó}{{\'o}}1 {ú}{{\'u}}1
           {à}{{\`a}}1 {â}{{\^a}}1 {ê}{{\^e}}1 {ô}{{\^o}}1
           {ã}{{\~a}}1 {õ}{{\~o}}1 {ç}{{\c c}}1
           {Á}{{\'A}}1 {É}{{\'E}}1 {Ç}{{\c C}}1 {Ã}{{\~A}}1
           {—}{{---}}1 {→}{{$\rightarrow$}}1 {⁻}{{$^-$}}1
}
\lstdefinelanguage{Verilog2}[]{Verilog}{morekeywords={localparam,signed}}
\lstdefinestyle{sh}{language=bash,keywordstyle=\color{black}}

\title{\textbf{Como adicionar um módulo ao Morphe}\\[4pt]
  \large Roteiro passo a passo, com o exemplo da soma de dois sinais}
\author{Antonio Nicassio Santos Lima\\
  \small Estágio no LabPS/UEFS --- supervisão: Prof.\ Armando S.\ Sanca}
\date{outubro de 2026}

\begin{document}
\maketitle

\begin{abstract}
Este roteiro mostra, do começo ao fim, como pôr uma operação nova no Morphe:
o circuito na FPGA, as memórias e os registradores que ligam esse circuito ao
processador da placa, o servidor em C que roda nesse processador e a janela no
aplicativo do aluno. O exemplo é deliberadamente simples --- a soma de dois
sinais, \(y[n] = a[n] + b[n]\) --- para que a atenção fique no caminho, que é o
mesmo para qualquer acelerador. As partes que são só procedimento aparecem por
cima; as que costumam dar errado (como o processador enxerga a memória da FPGA,
onde ficam os endereços, que arquivo vai em qual pasta) são explicadas em
detalhe. Sempre que o projeto já tem uma peça pronta, o roteiro mostra como
reaproveitá-la. Todo o código está na branch \arq{estagio/modulo-soma} do
repositório.
\end{abstract}

\tableofcontents

% =============================================================================
\section{Antes de começar}
% =============================================================================

\subsection{O que o módulo faz}

Recebe dois sinais de até 1024 amostras, \(a[n]\) e \(b[n]\), e devolve
\(y[n] = a[n] + b[n]\), amostra a amostra. Os números estão em
\textbf{Q15.16}, o formato de todo o Morphe: inteiros de 32 bits com sinal,
dos quais os 16 de baixo são a parte fracionária (o valor real é o inteiro
dividido por \(2^{16} = 65\,536\)). Somar dois Q15.16 dá outro Q15.16, mas a
soma pode passar de 32 bits; nesse caso o resultado \textbf{satura} no maior
(ou no menor) valor representável, em vez de ``dar a volta'' e virar um número
de sinal trocado.

\subsection{As camadas}

Uma operação atravessa quatro camadas, e o módulo novo precisa de uma peça em
cada uma (Figura~\ref{fig:camadas}):

\begin{figure}[htbp]
\centering
\begin{tikzpicture}[
  caixa/.style={draw,rounded corners=2pt,minimum width=3.2cm,minimum height=1cm,
                align=center,font=\small},
  seta/.style={-{Stealth[length=2mm]},thick}]
  \node[caixa,fill=cliente!10,draw=cliente] (cli) {Aplicativo\\\footnotesize\arq{soma_window.py}};
  \node[caixa,fill=hps!10,draw=hps,right=1.4cm of cli] (srv) {Servidor C no HPS\\\footnotesize\arq{handle_soma()}};
  \node[caixa,fill=fpga!8,draw=fpga,right=1.4cm of srv] (pd) {Memórias e PIOs\\\footnotesize Platform Designer};
  \node[caixa,fill=fpga!15,draw=fpga,below=1.1cm of pd] (rtl) {Circuito\\\footnotesize\arq{soma.v}};
  \draw[seta] (cli) -- node[above,font=\scriptsize]{TCP} node[below,font=\scriptsize]{opcode 9} (srv);
  \draw[seta] (srv) -- node[above,font=\scriptsize]{/dev/mem} node[below,font=\scriptsize]{pontes} (pd);
  \draw[seta] (pd) -- node[right,font=\scriptsize]{porta s1, start/done} (rtl);
\end{tikzpicture}
\caption{O caminho de uma operação. O aplicativo fala com o servidor pela
rede; o servidor escreve e lê as memórias da FPGA pelas pontes HPS--FPGA; o
circuito lê e escreve as mesmas memórias pela outra porta.}
\label{fig:camadas}
\end{figure}

\subsection{A ordem de construção}

A ordem recomendada é \textbf{de baixo para cima}, depois de combinar o que
passa entre as camadas: primeiro o circuito, depois as memórias e o topo do
projeto, a compilação, o cabeçalho de endereços, o servidor e, por último, o
aplicativo. Cada camada só é escrita quando a de baixo já existe, e assim,
quando algo dá errado, a falha está na camada que acabou de entrar. Começar
pelo aplicativo parece mais natural, mas deixa tudo sem ter contra o que
testar até o fim.

\subsection{Os arquivos que serão tocados}

\begin{table}[htbp]
\centering\small
\caption{Tudo o que o módulo novo muda. ``Novo'' é arquivo criado; o resto é
edição de arquivo que já existe.}
\label{tab:arquivos}
\begin{tabularx}{\textwidth}{@{}lX@{}}
\toprule
\textbf{Arquivo} & \textbf{O que muda} \\
\midrule
\arq{Quartus/soma.v} & novo: o circuito \\
\arq{Quartus/soc_system.qsf} & uma linha: o \arq{soma.v} entra no projeto \\
\arq{Quartus/soc_system.qsys} & três memórias e dois PIOs (Platform Designer) \\
\arq{Quartus/ghrd_top.v} & fios, instância do \arq{soma} e portas do \arq{soc_system} \\
\arq{C/hps_0.h} & regenerado, nunca editado à mão \\
\arq{C/address_map_arm.h} & a janela de memória que o servidor mapeia cresce \\
\arq{C/morphe_config.h}, \arq{C/morphe_protocol.h} & o limite e o opcode novos \\
\arq{C/morphe_server.c} & o tratamento da operação e o despacho \\
\arq{Python/morphe_config.py}, \arq{Python/morphe_protocol.py} & espelho dos dois arquivos do C \\
\arq{Python/soma_window.py} & novo: a janela \\
\arq{Python/morphe_app.py} & o botão no menu \\
\bottomrule
\end{tabularx}
\end{table}

% =============================================================================
\section{O combinado entre as camadas}
% =============================================================================

Antes de escrever código, decida e anote quatro coisas. Elas aparecem em
vários arquivos, e um número diferente em um deles é o erro mais comum.

\paragraph{1. O opcode.} Cada operação tem um número no protocolo. Os usados
estão em \arq{C/morphe_protocol.h} (1 a 8); a soma é o \textbf{9}. O mesmo
número vai para o \arq{Python/morphe_protocol.py}.

\paragraph{2. O formato do pedido e da resposta.} Sempre que possível, copie o
formato de uma operação que já existe. A soma usa o da convolução e do FIR:
cabeçalho de 20 bytes com \texttt{n\_x} e \texttt{n\_h}, depois as
\texttt{n\_x} amostras de \(a\) e as \texttt{n\_h} de \(b\), todas
\texttt{int32} em Q15.16. A resposta traz \texttt{n\_x} amostras de \(y\).
Copiar o formato economiza código dos dois lados.

\paragraph{3. O limite.} 1024 amostras, como os outros blocos. Ele vira
\texttt{MORPHE\_SOMA\_N\_MAX} no C e \texttt{SOMA\_N\_MAX} no Python, e o
mesmo número é o parâmetro \texttt{N\_SAMPLES} do Verilog.

\paragraph{4. As memórias, os PIOs e os endereços.} O circuito precisa de
uma memória por vetor (\(a\), \(b\) e \(y\)) e de dois sinais de controle:
\texttt{start}, que o processador liga para disparar, e \texttt{done}, que o
circuito liga ao terminar. Os endereços saem dos buracos livres do mapa atual,
que está em \arq{C/hps_0.h} (seção~\ref{sec:enderecos}).

\begin{table}[htbp]
\centering\small
\caption{O que a soma acrescenta ao mapa de endereços.}
\label{tab:mapa}
\begin{tabular}{@{}llll@{}}
\toprule
\textbf{Nome} & \textbf{Ponte} & \textbf{Endereço} & \textbf{Tipo} \\
\midrule
\texttt{soma\_a} & \texttt{h2f\_axi\_master} & \texttt{0x40000} & memória, 4\,KiB (1024 × 32 bits) \\
\texttt{soma\_b} & \texttt{h2f\_axi\_master} & \texttt{0x41000} & memória, 4\,KiB \\
\texttt{soma\_y} & \texttt{h2f\_axi\_master} & \texttt{0x42000} & memória, 4\,KiB \\
\texttt{soma\_start} & \texttt{h2f\_lw\_axi\_master} & \texttt{0x130} & PIO de saída, 1 bit \\
\texttt{soma\_done} & \texttt{h2f\_lw\_axi\_master} & \texttt{0x140} & PIO de entrada, 1 bit \\
\bottomrule
\end{tabular}
\end{table}

\subsection{Como o processador enxerga a FPGA}
\label{sec:enderecos}

Esta é a parte que mais confunde, e vale entender de uma vez. O processador da
placa (o HPS, um ARM com Linux) e a FPGA conversam por \textbf{pontes}: faixas
de endereço físico do processador que, em vez de levarem à memória RAM, levam
a componentes dentro da FPGA. O Morphe usa duas:

\begin{description}[leftmargin=1.2em]
  \item[\texttt{h2f\_axi\_master}] a ponte larga e rápida, onde ficam as
    \textbf{memórias}. Começa no endereço físico \texttt{0xC0000000}
    (\texttt{FPGA\_ONCHIP\_BASE}, em \arq{C/address_map_arm.h}).
  \item[\texttt{h2f\_lw\_axi\_master}] a ponte leve (\emph{lightweight}),
    onde ficam os registradores de controle, os \textbf{PIOs}. Começa em
    \texttt{0xFF200000} (\texttt{LW\_BRIDGE\_BASE}).
\end{description}

O endereço que você escolhe no Platform Designer é um \textbf{deslocamento
dentro da ponte}. A memória \texttt{soma\_a} em \texttt{0x40000} da ponte
larga está, para o processador, no endereço físico
\(\texttt{0xC0000000} + \texttt{0x40000} = \texttt{0xC0040000}\); o PIO
\texttt{soma\_start} em \texttt{0x130} da ponte leve está em
\texttt{0xFF200130}.

\paragraph{Onde ver o mapa.} Três lugares mostram a mesma coisa:
\begin{itemize}
  \item a aba \emph{Address Map} do Platform Designer;
  \item o \arq{C/hps_0.h}, gerado do \arq{soc_system.sopcinfo} pelo
    \arq{gen_hps_header.py}: para cada componente, um \texttt{\_BASE} (o
    deslocamento) e um \texttt{\_SPAN} (o tamanho em bytes). É o que o servidor usa, e o jeito mais rápido de achar um buraco
    livre:
\begin{lstlisting}[style=sh]
grep -E "_(BASE|SPAN) " C/hps_0.h | sort -k3
\end{lstlisting}
  \item o próprio \arq{soc_system.sopcinfo}, que é XML e de onde os outros
    dois saem.
\end{itemize}

Antes da soma, a ponte larga estava ocupada até \texttt{0x3FFFF} (o
\texttt{adc\_buf}, de 128\,KiB, começa em \texttt{0x20000}); por isso as
memórias novas foram para \texttt{0x40000}. Na ponte leve, o último PIO era o
\texttt{adc\_contador}, em \texttt{0x120}, e os PIOs ocupam 16 bytes cada:
\texttt{0x130} e \texttt{0x140} estavam livres.

\paragraph{Como o servidor acessa esses endereços.} Um programa no Linux não
pode usar endereço físico direto. O servidor abre \arq{/dev/mem} e pede ao
kernel, com \texttt{mmap}, uma ``janela'' que corresponde a uma faixa física;
a partir daí, um ponteiro dentro da janela é um acesso à FPGA. O
\arq{morphe_server.c} já faz isso uma vez para cada ponte, na
\texttt{fpga\_init()}:

\begin{lstlisting}[language=C]
g_fpga_virt = mmap(NULL, MORPHE_ONCHIP_MMAP_SPAN, PROT_READ | PROT_WRITE,
                   MAP_SHARED, g_fd_mem, FPGA_ONCHIP_BASE);   /* ponte larga */
g_lw_virt   = mmap(NULL, LW_BRIDGE_SPAN, PROT_READ | PROT_WRITE,
                   MAP_SHARED, g_fd_mem, LW_BRIDGE_BASE);     /* ponte leve  */
\end{lstlisting}

Um componente novo é só mais um ponteiro: o início da janela mais o
deslocamento do \arq{hps_0.h}.

\begin{lstlisting}[language=C]
g_soma_a         = (int32_t *)((char *)g_fpga_virt + SOMA_A_BASE);
g_pio_soma_start = (volatile uint32_t *)((char *)g_lw_virt + SOMA_START_BASE);

g_soma_a[i] = valor;         /* escreve a palavra i da memoria soma_a */
*g_pio_soma_start = 1;       /* liga o start                          */
\end{lstlisting}

O \texttt{volatile} nos PIOs é obrigatório: ele diz ao compilador que o valor
pode mudar sozinho (o \texttt{done} é escrito pelo circuito) e que cada leitura
e escrita tem de acontecer de fato, na ordem do código.

\begin{atencao}
\textbf{Acessar um endereço da ponte onde não há componente trava a placa.}
No Cyclone~V, o barramento espera para sempre por uma resposta que não vem: a
placa some da rede e nem a serial responde. Isso acontece, por exemplo, quando
o servidor novo (que conhece a soma) roda numa placa programada com o
bitstream antigo (que não tem nada em \texttt{0x40000}). A
seção~\ref{sec:servidor} mostra a proteção que o projeto já tem contra isso.
\end{atencao}

% =============================================================================
\section{O gabarito}
% =============================================================================

Antes do circuito, escreva a mesma conta em Python, com os mesmos números
inteiros que o hardware vai usar. É ela que, no fim, diz se a FPGA acertou ---
e, num módulo em ponto fixo, a comparação certa é de \textbf{igualdade bit a
bit}, não de ``parecido''. A soma em Q15.16 com saturação é:

\begin{lstlisting}[language=Python]
def soma_referencia(a_q, b_q):
    """A conta do soma.v em NumPy: soma em 64 bits e satura na faixa de 32."""
    s = np.asarray(a_q, dtype=np.int64) + np.asarray(b_q, dtype=np.int64)
    return np.clip(s, dsp.Q_MIN, dsp.Q_MAX)
\end{lstlisting}

\begin{pronto}
A conversão de \texttt{float} para Q15.16 e a volta já existem em
\arq{Python/dsp_core.py}: \texttt{float\_to\_q1516()},
\texttt{q1516\_to\_float()}, e os limites \texttt{Q\_MIN} e \texttt{Q\_MAX}.
Não escreva outras.
\end{pronto}

% =============================================================================
\section{O circuito}
% =============================================================================

\subsection{Onde o arquivo fica}

O Verilog novo vai na \textbf{raiz da pasta} \arq{Quartus/}, ao lado do
\arq{conv1d.v} e do \arq{iir_cascade.v}: \arq{Quartus/soma.v}. E entra no
projeto por uma linha no \arq{Quartus/soc_system.qsf} --- ou pelo menu do
Quartus \emph{Project \(\rightarrow\) Add/Remove Files in Project}, que
escreve a mesma linha:

\begin{lstlisting}[style=sh]
set_global_assignment -name VERILOG_FILE soma.v
\end{lstlisting}

\begin{atencao}
Não ponha o Verilog numa subpasta nem o deixe fora do \arq{.qsf}. Alguns
arquivos do projeto (os \arq{iir_*.v}) não estão listados no \arq{.qsf}: o
Quartus os acha só por estarem na pasta do projeto, e fora dela a compilação
falha. Listar o arquivo novo no \arq{.qsf} é o jeito explícito, que não
depende disso.
\end{atencao}

\subsection{O contrato com o resto}

Todo acelerador do Morphe segue o mesmo protocolo de controle, e o novo deve
seguir também, porque o servidor já sabe conduzi-lo:

\begin{enumerate}
  \item o processador escreve as entradas nas memórias e liga
    \texttt{start}; o circuito dispara na \textbf{borda de subida} (assim um
    \texttt{start} que ficou ligado não dispara de novo);
  \item o circuito processa \textbf{sempre} \texttt{N\_SAMPLES} amostras ---
    o servidor completa com zeros o que o aluno não mandou --- e liga
    \texttt{done};
  \item o processador lê a saída e desliga \texttt{start}; com
    \texttt{start} em zero, o circuito desliga \texttt{done} e volta ao
    repouso.
\end{enumerate}

\subsection{Ler e escrever as memórias: os controladores prontos}

\begin{pronto}
O projeto tem dois módulos que fazem o protocolo de uma memória do Platform
Designer: \arq{memory_read_controller.v} e \arq{memory_write_controller.v}. O
\arq{conv1d.v} e o \arq{iir_cascade.v} usam os dois. Você dá o endereço e um
pulso de um ciclo em \texttt{start\_read} (ou \texttt{start\_write}); o
controlador cuida de \texttt{chipselect}, \texttt{clken} e da latência da
memória, e devolve um pulso em \texttt{read\_done} (ou
\texttt{write\_done}) com o dado pronto em \texttt{data\_out}.
\end{pronto}

A soma usa três: dois de leitura, para \(a\) e \(b\) em paralelo, e um de
escrita, para \(y\). Cada um fala com uma memória:

\begin{lstlisting}[language=Verilog2]
memory_read_controller #(
    .DATA_WIDTH         (DATA_WIDTH),
    .MEM_ADDRESS_N_BITS (ADDR_N_BITS)
) u_rd_a (
    .clk (clk), .reset_n (reset_n),
    .start_read    (rd_start),          // pulso: le a[n]
    .sram_readdata (a_sram_readdata),   // vem da memoria soma_a
    .data_addr     (n),
    .mem_write     (a_sram_write),      // vai para a memoria (sempre 0)
    .data_out      (a_data),            // a[n], valido com a_done
    .mem_address   (a_sram_address),
    .read_done     (a_done),
    .clken         (a_sram_clken),
    .chipselect    (a_sram_chipselect)
);
\end{lstlisting}

\subsection{A máquina de estados}

Com os controladores, o módulo vira uma máquina de seis estados: repouso, pede
a leitura, espera, pede a escrita, espera, terminou. A conta em si é uma linha
--- a soma em 33 bits e a saturação:

\begin{lstlisting}[language=Verilog2]
wire signed [DATA_WIDTH:0] soma_larga =
    $signed({a_data[DATA_WIDTH-1], a_data}) + $signed({b_data[DATA_WIDTH-1], b_data});

wire [DATA_WIDTH-1:0] soma_sat =
    (soma_larga > MAXV) ? MAXV[DATA_WIDTH-1:0] :
    (soma_larga < MINV) ? MINV[DATA_WIDTH-1:0] :
                          soma_larga[DATA_WIDTH-1:0];
\end{lstlisting}

\begin{lstlisting}[language=Verilog2]
S_IDLE: begin
    done <= 1'b0;
    if (start_pulse) begin n <= 0; state <= S_READ; end
end
S_READ:      begin rd_start <= 1'b1; state <= S_WAIT_READ; end
S_WAIT_READ: if (a_done && b_done) begin wr_data <= soma_sat; state <= S_WRITE; end
S_WRITE:     begin wr_start <= 1'b1; state <= S_WAIT_WRITE; end
S_WAIT_WRITE:
    if (wr_done) begin
        if (n == N_SAMPLES - 1) state <= S_DONE;
        else begin n <= n + 1'b1; state <= S_READ; end
    end
S_DONE: begin
    done <= 1'b1;
    if (!start) state <= S_IDLE;      // espera o HPS desligar o start
end
\end{lstlisting}

O arquivo completo está no apêndice~\ref{ap:somav}. Medido em simulação: uma
amostra custa 12 ciclos, e as 1024 levam 12\,289 ciclos, cerca de 0,25\,ms a
50\,MHz.

% =============================================================================
\section{As memórias, os PIOs e o topo do projeto}
% =============================================================================

\subsection{No Platform Designer}

Abra o Quartus 20.1, o projeto \arq{Quartus/soc_system.qpf} e, em
\emph{Tools \(\rightarrow\) Platform Designer}, o \arq{soc_system.qsys}.

\paragraph{Cada memória} (\texttt{soma\_a}, \texttt{soma\_b}, \texttt{soma\_y}):
\begin{enumerate}
  \item no \emph{IP Catalog}, \emph{Basic Functions \(\rightarrow\) On Chip
    Memory \(\rightarrow\) On-Chip Memory (RAM or ROM) Intel FPGA IP};
  \item tipo \emph{RAM}, \textbf{\emph{Dual-port access} marcado}, largura
    de dados 32, tamanho total \textbf{4096 bytes} (1024 palavras), tipo de
    bloco M10K, latência 1 nas duas portas; \emph{Finish};
  \item renomeie a instância (botão direito, \emph{Rename}) para
    \texttt{soma\_a};
  \item ligue \texttt{clk1} e \texttt{clk2} em \texttt{clk\_0.clk},
    \texttt{reset1} e \texttt{reset2} em \texttt{clk\_0.clk\_reset};
  \item ligue a porta \textbf{\texttt{s2}} em
    \texttt{hps\_0.h2f\_axi\_master}: é por ela que o processador escreve e
    lê;
  \item na coluna \emph{Export}, exporte a porta \textbf{\texttt{s1}} com o
    nome \texttt{soma\_a}: é por ela que o \arq{soma.v} lê.
\end{enumerate}

\paragraph{Cada PIO} (\texttt{soma\_start}, \texttt{soma\_done}):
\begin{enumerate}
  \item \emph{Processors and Peripherals \(\rightarrow\) Peripherals
    \(\rightarrow\) PIO (Parallel I/O) Intel FPGA IP};
  \item largura 1; direção \emph{Output} para o \texttt{soma\_start} (o
    processador escreve) e \emph{Input} para o \texttt{soma\_done} (o
    processador lê);
  \item \texttt{clk} e \texttt{reset} como nas memórias; \texttt{s1} em
    \texttt{hps\_0.h2f\_lw\_axi\_master} (a ponte \textbf{leve});
  \item exporte \texttt{external\_connection} com o nome do PIO.
\end{enumerate}

\paragraph{Os endereços.} Na coluna \emph{Base}, digite os da
Tabela~\ref{tab:mapa} e trave-os (o cadeado ao lado).

\begin{atencao}
\textbf{Não use \emph{System \(\rightarrow\) Assign Base Addresses}.} Ele
redistribui os endereços de \emph{todos} os componentes, e então o servidor, o
\arq{hps_0.h} versionado e o bitstream das placas deixam de concordar entre si.
Escolha os endereços novos à mão, nos buracos livres, e não mova os antigos.
\end{atencao}

\paragraph{Sem a interface gráfica.} O \arq{.qsys} é XML. Os componentes da
soma foram criados na branch copiando os blocos de um módulo que já existia
(\texttt{iir\_xn} para as memórias, \texttt{iir\_start} e
\texttt{iir\_done} para os PIOs) e trocando nomes e endereços. Funciona e é
reproduzível, mas exige cuidado: cada componente aparece em quatro lugares do
arquivo (o \texttt{element} com o endereço, a \texttt{interface} exportada, o
\texttt{module} com os parâmetros e as \texttt{connection} de relógio, reset e
barramento), e o nome aparece também colado, como em
\texttt{\$\$\{FILENAME\}\_iir\_xn} (o arquivo de inicialização da memória),
que tem de mudar junto. Na dúvida, use a interface gráfica.

\subsection{Gerar o sistema}

Depois de salvar, o sistema precisa ser \textbf{gerado}: é a geração que
escreve o Verilog do \arq{soc_system} (com as portas novas) e o
\arq{soc_system.sopcinfo} (com os endereços novos). Na estação, de dentro de
\arq{Quartus/}:

\begin{lstlisting}[style=sh]
qsys-generate soc_system.qsys --synthesis=VERILOG --family="Cyclone V" --part=5CSEMA5F31C6
\end{lstlisting}

\begin{atencao}
Gere sempre com o Quartus \textbf{20.1}, a versão do projeto. Uma versão mais
nova ``atualiza'' os componentes e muda arquivos que nada têm a ver com o
módulo novo. Se o comando não for encontrado, o terminal está em \texttt{ksh}:
digite \texttt{bash} antes.
\end{atencao}

\subsection{Onde ver os nomes das portas geradas}

O \arq{soc_system} gerado ganha uma porta para cada sinal exportado, com o
nome \emph{exportação}\texttt{\_}\emph{sinal}. A lista está no começo de
\arq{Quartus/soc_system/synthesis/soc_system.v}:

\begin{lstlisting}[style=sh]
grep -n "soma" Quartus/soc_system/synthesis/soc_system.v | head -30
\end{lstlisting}

Para uma memória exportada como \texttt{soma\_a} saem \texttt{soma\_a\_address},
\texttt{\_clken}, \texttt{\_chipselect}, \texttt{\_write}, \texttt{\_readdata},
\texttt{\_writedata} e \texttt{\_byteenable}; para um PIO,
\texttt{soma\_start\_export}. Não use o \arq{soc_system/soc_system_inst.v}: ele
só é atualizado pela interface gráfica e, no repositório, está desatualizado.

\subsection{O \texttt{ghrd\_top.v}}

O \arq{Quartus/ghrd_top.v} é o topo do projeto: é ali que o \arq{soc_system}
e os aceleradores são instanciados e ligados. A soma entra em quatro lugares,
cada um ao lado do bloco equivalente do IIR, que serve de modelo:

\begin{enumerate}
  \item \textbf{parâmetros}, junto dos outros \texttt{localparam}:
\begin{lstlisting}[language=Verilog2]
localparam SOMA_DATA_WIDTH = 32;
localparam SOMA_ADDR_WIDTH = 10;        // 1024 amostras
localparam SOMA_N_SAMPLES  = 1024;
\end{lstlisting}
  \item \textbf{fios}: dois para os PIOs e sete por memória;
\begin{lstlisting}[language=Verilog2]
wire                       soma_start;      // HPS -> FPGA
wire                       soma_done;       // FPGA -> HPS
wire [SOMA_ADDR_WIDTH-1:0] soma_a_address;
wire                       soma_a_clken, soma_a_chipselect, soma_a_write;
wire [SOMA_DATA_WIDTH-1:0] soma_a_readdata, soma_a_writedata;
wire [3:0]                 soma_a_byteenable = 4'b1111;
// ... o mesmo para soma_b e soma_y
\end{lstlisting}
  \item \textbf{a instância do \texttt{soma}}, ligada a esses fios, com o
    relógio de 50\,MHz e o reset do HPS:
\begin{lstlisting}[language=Verilog2]
soma #(.DATA_WIDTH(SOMA_DATA_WIDTH), .N_SAMPLES(SOMA_N_SAMPLES),
       .ADDR_N_BITS(SOMA_ADDR_WIDTH)) soma_inst (
    .clk (CLOCK_50), .reset_n (hps_fpga_reset_n),
    .start (soma_start), .done (soma_done),
    .a_sram_readdata (soma_a_readdata), .a_sram_address (soma_a_address),
    .a_sram_chipselect (soma_a_chipselect), .a_sram_clken (soma_a_clken),
    .a_sram_write (soma_a_write),
    // ... b igual a a; y com .y_sram_wdata (soma_y_writedata)
);
\end{lstlisting}
  \item \textbf{as portas novas na instância \texttt{u0} do
    \texttt{soc\_system}}, com os nomes da subseção anterior:
\begin{lstlisting}[language=Verilog2]
.soma_start_export  (soma_start),
.soma_done_export   (soma_done),
.soma_a_address     (soma_a_address),
.soma_a_clken       (soma_a_clken),
// ... todos os sinais das tres memorias
\end{lstlisting}
\end{enumerate}

Repare na direção: \texttt{soma\_a} e \texttt{soma\_b} são lidas pelo
circuito (\texttt{readdata} entra no \arq{soma.v}); \texttt{soma\_y} é escrita
(\texttt{writedata} sai dele). O \texttt{byteenable} fica em
\texttt{4'b1111}: o circuito sempre escreve a palavra inteira.

% =============================================================================
\section{Compilar}
% =============================================================================

Na estação, no clone de desenvolvimento (não no \arq{/opt/morphe} da turma),
depois do \texttt{qsys-generate}:

\begin{lstlisting}[style=sh]
quartus_sh --flow compile soc_system
\end{lstlisting}

Leva cerca de 30 minutos. No fim, confira dois resumos:

\begin{lstlisting}[style=sh]
cat output_files/soc_system.fit.summary
grep -A2 "clock_50_1" output_files/soc_system.sta.summary | head -12
\end{lstlisting}

\begin{confira}
\begin{itemize}[leftmargin=1em]
  \item \emph{Full Compilation was successful};
  \item no \arq{.fit.summary}, mais memória que antes (as três memórias novas
    usam cerca de 12 blocos M10K) e um pouco mais de lógica;
  \item no \arq{.sta.summary}, nenhuma folga (\emph{Slack}) negativa no relógio
    \texttt{clock\_50\_1}. Folga negativa quer dizer que o circuito não fecha a
    50\,MHz; o \arq{Quartus/relatorio_timing.tcl} mostra o caminho culpado.
\end{itemize}
\end{confira}

O bitstream sai em \arq{output_files/soc_system_time_limited.sof}. O
``\emph{time limited}'' vem do núcleo de FFT, que é IP licenciada: por isso o
\arq{morphe-up.sh} mantém o \emph{tether} da licença rodando enquanto a placa
é usada.

% =============================================================================
\section{O cabeçalho de endereços}
% =============================================================================

O servidor não tem endereço escrito à mão: ele os lê do \arq{C/hps_0.h}, que é
\textbf{gerado} a partir do \arq{.sopcinfo} produzido no passo anterior.

\begin{lstlisting}[style=sh]
python3 Quartus/gen_hps_header.py
grep -E "SOMA_|SYSID_QSYS_TIMESTAMP" C/hps_0.h
\end{lstlisting}

\begin{confira}
Têm de aparecer \texttt{SOMA\_A\_BASE 0x40000}, \texttt{SOMA\_B\_BASE
0x41000}, \texttt{SOMA\_Y\_BASE 0x42000} (cada uma com \texttt{\_SPAN 4096}),
\texttt{SOMA\_START\_BASE 0x130}, \texttt{SOMA\_DONE\_BASE 0x140} e um
\texttt{SYSID\_QSYS\_TIMESTAMP} novo. Depois, \texttt{python3
Quartus/gen\_hps\_header.py -{}-check} confere que o cabeçalho e o
\arq{.sopcinfo} concordam.
\end{confira}

\begin{atencao}
Nunca edite o \arq{hps_0.h} à mão. Um cabeçalho que descreve hardware que não
está na FPGA leva o servidor a acessar endereço sem componente --- e trava a
placa (seção~\ref{sec:enderecos}).
\end{atencao}

% =============================================================================
\section{O servidor C}
\label{sec:servidor}
% =============================================================================

Tudo em \arq{C/}. O molde é a operação FIR (\texttt{handle\_fir}), que tem o
mesmo formato de pedido.

\subsection{Os limites e o opcode}

\begin{lstlisting}[language=C]
/* morphe_config.h */
#define MORPHE_SOMA_N_MAX       1024

/* morphe_protocol.h */
#define MORPHE_OP_SOMA     9U   /* y = a + b (soma.v) */
\end{lstlisting}

\subsection{A janela de memória}

O servidor mapeia uma janela só para a ponte larga, do começo até o fim da
última memória. Esse tamanho é calculado sozinho a partir do \arq{hps_0.h}
(\texttt{MORPHE\_ONCHIP\_MMAP\_SPAN}), mas não pode passar do
\texttt{FPGA\_ONCHIP\_SPAN} de \arq{address_map_arm.h}. Como as memórias da
soma passam de \texttt{0x3FFFF}, esse limite cresce:

\begin{lstlisting}[language=C]
/* address_map_arm.h */
#define FPGA_ONCHIP_SPAN       0x0004FFFF    /* era 0x0003FFFF */
\end{lstlisting}

Se esquecer, o compilador avisa: o \texttt{\_Static\_assert} que compara os
dois falha. E a memória nova tem de entrar no cálculo do tamanho da janela:

\begin{lstlisting}[language=C]
#define _MORPHE_END_SOMA \
    _MORPHE_MAX3( \
        _MORPHE_END(SOMA_A_BASE, SOMA_A_SPAN), \
        _MORPHE_END(SOMA_B_BASE, SOMA_B_SPAN), \
        _MORPHE_END(SOMA_Y_BASE, SOMA_Y_SPAN))
/* e _MORPHE_END_SOMA entra no MORPHE_ONCHIP_MMAP_SPAN */
\end{lstlisting}

\subsection{Compilar a soma só quando o hardware a tem}

\begin{pronto}
O ADC resolveu, antes, o problema de um servidor novo numa placa com
bitstream velho. Duas peças, que a soma reaproveita:
\begin{itemize}[leftmargin=1em]
  \item \textbf{compilação condicional}: o código só entra se o
    \arq{hps_0.h} tem os símbolos do componente --- ou seja, se o bitstream
    que gerou o cabeçalho tem o componente;
  \item \textbf{o \emph{sysid}}: um registrador da FPGA guarda o instante em
    que o sistema foi gerado (\texttt{SYSID\_QSYS\_TIMESTAMP}). O servidor só
    toca no componente se o número gravado na FPGA é o mesmo do seu
    \arq{hps_0.h}. Caso contrário, recusa o pedido com uma mensagem, sem
    acessar a FPGA.
\end{itemize}
\end{pronto}

\begin{lstlisting}[language=C]
#if defined(SOMA_A_BASE) && defined(SOMA_B_BASE) && defined(SOMA_Y_BASE) && \
    defined(SOMA_START_BASE) && defined(SOMA_DONE_BASE) && MORPHE_TEM_ADC
#define MORPHE_TEM_SOMA 1
#else
#define MORPHE_TEM_SOMA 0
#endif
\end{lstlisting}

(O \texttt{MORPHE\_TEM\_ADC} entra na condição porque é ele que garante que
o servidor lê o \emph{sysid}.)

\subsection{Os ponteiros}

Declarados junto dos outros e preenchidos na \texttt{fpga\_init()}, como na
seção~\ref{sec:enderecos}:

\begin{lstlisting}[language=C]
#if MORPHE_TEM_SOMA
static int32_t           *g_soma_a, *g_soma_b, *g_soma_y;
static volatile uint32_t *g_pio_soma_start, *g_pio_soma_done;
#endif

/* em fpga_init(): */
g_soma_a         = (int32_t *)((char *)g_fpga_virt + SOMA_A_BASE);
g_soma_b         = (int32_t *)((char *)g_fpga_virt + SOMA_B_BASE);
g_soma_y         = (int32_t *)((char *)g_fpga_virt + SOMA_Y_BASE);
g_pio_soma_start = (volatile uint32_t *)((char *)g_lw_virt + SOMA_START_BASE);
g_pio_soma_done  = (volatile uint32_t *)((char *)g_lw_virt + SOMA_DONE_BASE);
\end{lstlisting}

\begin{atencao}
A \texttt{fpga\_init()} só calcula ponteiros, \textbf{nunca acessa a FPGA}:
o servidor sobe no boot da placa, antes de o \arq{morphe-up.sh} programar o
bitstream do Morphe, e um acesso nesse momento trava a placa. O primeiro
acesso (zerar o \texttt{start}) fica na \texttt{fpga\_preparada()}, protegido
pelo \emph{sysid}:
\begin{lstlisting}[language=C]
if (g_sysid[1] == SYSID_QSYS_TIMESTAMP) *g_pio_soma_start = 0;
\end{lstlisting}
\end{atencao}

\subsection{O tratamento da operação}

A função \texttt{handle\_soma()} segue o \texttt{handle\_fir()} passo a passo
(o código completo está no apêndice~\ref{ap:handle}):

\begin{enumerate}
  \item recusa se a FPGA não foi preparada (\texttt{fpga\_preparada()}) ou
    se o \emph{sysid} não confere;
  \item confere os tamanhos (\(a\) e \(b\) iguais, de 1 a
    \texttt{MORPHE\_SOMA\_N\_MAX}) e responde com erro se não;
  \item recebe o payload (\texttt{recv\_exact}) e o converte
    (\texttt{decode\_samples\_to\_i32}, que já trata \texttt{int32} e
    \texttt{float32});
  \item escreve \(a\) e \(b\) nas memórias, completando com zeros até 1024;
  \item dispara e espera:
\begin{lstlisting}[language=C]
*g_pio_soma_start = 0;
usleep(1);
*g_pio_soma_start = 1;                    /* borda de subida */
if (wait_done(g_pio_soma_done, FPGA_DONE_TIMEOUT_MS) < 0) { /* timeout */ }
*g_pio_soma_start = 0;
\end{lstlisting}
  \item lê \(y\) e responde (\texttt{build\_resp\_header} e
    \texttt{send\_all}).
\end{enumerate}

\begin{pronto}
Tudo o que a função usa já existe no servidor: \texttt{recv\_exact},
\texttt{send\_all}, \texttt{send\_error}, \texttt{decode\_samples\_to\_i32},
\texttt{build\_resp\_header}, \texttt{wait\_done} (espera o \texttt{done} com
prazo) e \texttt{save\_debug\_bundle\_conv}, que grava na placa um
\arq{.mrph} de cada operação para depuração --- o da soma usa o mesmo formato
da convolução, com \(a\) no lugar de \(x\) e \(b\) no de \(h\).
\end{pronto}

\subsection{Ligar a função ao resto}

Três linhas fecham o servidor:

\begin{lstlisting}[language=C]
/* no despacho, serve_connection(): */
case MORPHE_OP_SOMA: handle_soma(sock, dtype, n_x, n_h); break;

/* em nome_op(), o nome que aparece na porta de estado e no painel: */
case MORPHE_OP_SOMA: return "soma";

/* no PING, para o cliente saber se a placa tem a soma: */
"soma_n_max=%d\n"   ->   soma_ok ? MORPHE_SOMA_N_MAX : 0
\end{lstlisting}

O servidor não precisa ser compilado à mão: o \arq{morphe-up.sh} percebe que
as fontes mudaram, envia-as à placa e recompila lá.

% =============================================================================
\section{Na placa}
% =============================================================================

Do clone de desenvolvimento, na estação, com o cabo da placa escolhida para o
teste (de preferência fora do horário de aula, porque a placa é reprogramada):

\begin{lstlisting}[style=sh]
./morphe-up.sh --cable 'DE-SoC [1-3]' --deploy
\end{lstlisting}

O script grava o \arq{.sof} \emph{deste clone}, mantém o \emph{tether}, envia
e recompila o servidor e roda os quatro testes do \arq{morphe_ping.py}.

\begin{confira}
\begin{itemize}[leftmargin=1em]
  \item os quatro testes do \arq{morphe_ping.py} passam: as operações antigas
    continuam funcionando com o bitstream novo;
  \item no aplicativo, a janela da soma devolve o resultado e a barra de
    status diz que as amostras são \textbf{iguais, bit a bit}, à referência;
  \item num pedido que satura (por exemplo, \(a = b = 30\,000\)), o
    resultado é \(32\,767{,}99998\), e não um número negativo.
\end{itemize}
\end{confira}

Para devolver a placa à turma, rode o \arq{morphe-up.sh} de
\arq{/opt/morphe} com o mesmo cabo: ele grava de volta o bitstream oficial.

% =============================================================================
\section{O aplicativo}
% =============================================================================

Tudo em \arq{Python/}.

\paragraph{Espelho do C.} O limite e o opcode, com os mesmos números:
\begin{lstlisting}[language=Python]
# morphe_config.py
SOMA_N_MAX: int = 1024

# morphe_protocol.py
OP_SOMA = 9
\end{lstlisting}

\paragraph{O pedido e a resposta.} Em \arq{morphe_protocol.py}, ao lado dos
da convolução, que têm o mesmo formato:
\begin{lstlisting}[language=Python]
def build_soma_request(a_q, b_q, dtype_code=DTYPE_INT32) -> bytes:
    hdr = struct.pack(">IHHHHII", MAGIC_REQ, VERSION, OP_SOMA, dtype_code, 0,
                      len(a_q), len(b_q))
    return hdr + pack_samples(a_q, dtype_code) + pack_samples(b_q, dtype_code)

def decode_soma_response(resp) -> np.ndarray:
    if not resp.ok:
        raise RuntimeError(resp.payload.decode("utf-8", "replace"))
    be = _np_dtype_be(resp.dtype_code)
    return np.frombuffer(resp.payload, dtype=be, count=resp.n_out).astype(np.int64)
\end{lstlisting}
O \texttt{TcpClient.request()} não muda: ele já sabe ler uma resposta de
\texttt{n\_out} palavras de 32 bits.

\paragraph{A janela.} \arq{soma_window.py} (completa no
apêndice~\ref{ap:janela}) tem cerca de 170 linhas, porque quase tudo já
existe:

\begin{pronto}
\begin{itemize}[leftmargin=1em]
  \item \texttt{SignalPanel} (\arq{signal_panel.py}): o painel de sinal de
    todas as janelas, com os tipos gerados e a leitura de arquivo;
  \item \texttt{\_stem\_signal} (\arq{conv_window.py}): desenha um sinal
    como na janela de convolução;
  \item \texttt{float\_to\_q1516}, \texttt{q1516\_to\_float},
    \texttt{pad\_zeros\_to} (\arq{dsp_core.py});
  \item \texttt{self.master.tcp\_panel.make\_client()}: o cliente já
    apontado para a placa escolhida no painel ``Placas do laboratório'';
  \item \texttt{theme} (\arq{morphe_theme.py}) e \texttt{PlotToolbar}: a
    aparência e a barra dos gráficos.
\end{itemize}
\end{pronto}

O único cuidado próprio de uma janela é \textbf{não esperar a rede na
\emph{thread} da interface}: o pedido vai numa \emph{thread} separada, e o
resultado volta à interface com \texttt{self.after(0, ...)}, porque o Tk só
pode ser mexido pela \emph{thread} dele.

\begin{lstlisting}[language=Python]
def worker():
    try:
        resp = client.request(build_soma_request(a_q, b_q, DTYPE_INT32))
        y_q = decode_soma_response(resp)
        self.after(0, lambda: self._on_resultado(y_q, a_q, b_q, n0))
    except Exception as e:
        self.after(0, lambda err=e: self._on_erro(err))

threading.Thread(target=worker, daemon=True).start()
\end{lstlisting}

\paragraph{O botão.} Em \arq{morphe_app.py}, um import, uma linha na lista
de botões e um método:
\begin{lstlisting}[language=Python]
from soma_window import SomaWindow
...
("Soma  (a + b → FPGA, exemplo do roteiro)", self._open_soma),
...
def _open_soma(self):
    SomaWindow(self)
\end{lstlisting}

% =============================================================================
\section{Fechar}
% =============================================================================

Com tudo validado na placa:

\begin{enumerate}
  \item \textbf{proveniência}: na estação, \texttt{./gera\_proveniencia.sh
    "como foi validado"} registra os \emph{hashes} do conjunto inseparável ---
    RTL, bitstream, \arq{hps_0.h} e servidor;
  \item \textbf{commit} do \arq{.sof}, do \arq{.sopcinfo}, do \arq{hps_0.h}
    e da pasta \arq{soc_system/} gerada (sem \arq{db/}), junto com as fontes;
  \item \textbf{mesclar} a branch na principal e fazer o \emph{push};
  \item \textbf{atualizar} a estação e os computadores:
    \texttt{/opt/morphe/atualizar.sh} em cada um;
  \item \textbf{preparar as placas} com o bitstream novo:
    \texttt{./morphe-up.sh -{}-todas}, de \arq{/opt/morphe}.
\end{enumerate}

\begin{atencao}
Os passos 4 e 5 andam juntos. O servidor novo numa placa com o bitstream
antigo recusa a soma (por causa do \emph{sysid}), mas não trava; o contrário
--- bitstream novo, servidor antigo --- funciona para tudo menos a soma.
\end{atencao}

% =============================================================================
\section{Armadilhas conhecidas}
% =============================================================================

\begin{description}[leftmargin=1em,style=nextline]
  \item[Verilog fora da pasta do projeto ou fora do \arq{.qsf}]
    A compilação não acha o módulo. Arquivo novo na raiz de \arq{Quartus/} e
    listado no \arq{.qsf}.
  \item[\texttt{Assign Base Addresses} no Platform Designer]
    Muda os endereços de tudo. Escolha os novos à mão e não mova os antigos.
  \item[Gerar ou compilar com outra versão do Quartus]
    Componentes ``atualizados'' e arquivos mudados sem relação com o módulo.
    Sempre a 20.1.
  \item[\arq{hps_0.h} diferente do bitstream]
    O servidor acessa endereço sem componente e a placa trava. Gere o
    cabeçalho com o \arq{gen_hps_header.py} e confira com \opt{check}.
  \item[Esquecer o \texttt{FPGA\_ONCHIP\_SPAN}]
    O \texttt{\_Static\_assert} falha na compilação do servidor. Aumente-o
    quando uma memória passar do fim da janela.
  \item[Acessar a FPGA na \texttt{fpga\_init()}]
    O servidor sobe no boot, antes do bitstream do Morphe. O primeiro acesso
    vai para a \texttt{fpga\_preparada()}.
  \item[O \emph{tether} da licença]
    O bitstream tem a FFT licenciada: sem o \emph{tether} rodando na estação,
    a placa para em cerca de uma hora. O \arq{morphe-up.sh} cuida disso; o
    \texttt{quartus\_pgm} à mão, não.
  \item[\texttt{sudo ./morphe-up.sh}]
    Usa o Quartus da conta \texttt{root} e deixa o \emph{tether} como
    processo dela. Rode sempre sem \texttt{sudo}.
  \item[Fins de linha do \arq{soc_system.qsf}]
    O arquivo tem fins de linha misturados (LF e CRLF). Editado num editor do
    Windows, o \emph{diff} pode mostrar o arquivo inteiro como alterado. Edite
    no Linux, ou pelo Quartus.
  \item[Terminal em \texttt{ksh}]
    \texttt{qsys-generate: not found}. Digite \texttt{bash} antes.
\end{description}

% =============================================================================
\appendix
\section{O \texttt{soma.v} completo}
\label{ap:somav}
% =============================================================================

\begin{lstlisting}[language=Verilog2]
`timescale 1 ns / 1 ps

module soma #(
    parameter DATA_WIDTH  = 32,
    parameter N_SAMPLES   = 1024,
    parameter ADDR_N_BITS = 10
)(
    input  wire                     clk,
    input  wire                     reset_n,
    input  wire                     start,
    output reg                      done,

    // --- memoria da entrada a[n] (so leitura) ---
    input  wire [DATA_WIDTH-1:0]    a_sram_readdata,
    output wire [ADDR_N_BITS-1:0]   a_sram_address,
    output wire                     a_sram_chipselect,
    output wire                     a_sram_clken,
    output wire                     a_sram_write,

    // --- memoria da entrada b[n] (so leitura) ---
    input  wire [DATA_WIDTH-1:0]    b_sram_readdata,
    output wire [ADDR_N_BITS-1:0]   b_sram_address,
    output wire                     b_sram_chipselect,
    output wire                     b_sram_clken,
    output wire                     b_sram_write,

    // --- memoria da saida y[n] (so escrita) ---
    output wire [ADDR_N_BITS-1:0]   y_sram_address,
    output wire [DATA_WIDTH-1:0]    y_sram_wdata,
    output wire                     y_sram_chipselect,
    output wire                     y_sram_clken,
    output wire                     y_sram_write
);

    localparam [2:0] S_IDLE       = 3'd0;
    localparam [2:0] S_READ       = 3'd1;  // pede a[n] e b[n]
    localparam [2:0] S_WAIT_READ  = 3'd2;
    localparam [2:0] S_WRITE      = 3'd3;  // grava y[n]
    localparam [2:0] S_WAIT_WRITE = 3'd4;
    localparam [2:0] S_DONE       = 3'd5;

    reg [2:0] state;
    reg [ADDR_N_BITS-1:0] n;    // amostra corrente

    // start dispara na borda de subida
    reg start_d;
    wire start_pulse = start & ~start_d;
    always @(posedge clk or negedge reset_n) begin
        if (!reset_n) start_d <= 1'b0;
        else          start_d <= start;
    end

    // --- controle das memorias ---
    reg                     rd_start;
    wire [DATA_WIDTH-1:0]   a_data, b_data;
    wire                    a_done, b_done;

    reg                     wr_start;
    reg  [DATA_WIDTH-1:0]   wr_data;
    wire                    wr_done;

    // --- a conta: 33 bits, depois satura em 32 ---
    localparam signed [DATA_WIDTH:0] MAXV = {2'b00, {(DATA_WIDTH-1){1'b1}}};  //  2^31 - 1
    localparam signed [DATA_WIDTH:0] MINV = {2'b11, {(DATA_WIDTH-1){1'b0}}};  // -2^31

    wire signed [DATA_WIDTH:0] soma_larga =
        $signed({a_data[DATA_WIDTH-1], a_data}) + $signed({b_data[DATA_WIDTH-1], b_data});

    wire [DATA_WIDTH-1:0] soma_sat =
        (soma_larga > MAXV) ? MAXV[DATA_WIDTH-1:0] :
        (soma_larga < MINV) ? MINV[DATA_WIDTH-1:0] :
                              soma_larga[DATA_WIDTH-1:0];

    always @(posedge clk or negedge reset_n) begin
        if (!reset_n) begin
            state    <= S_IDLE;
            done     <= 1'b0;
            n        <= 0;
            rd_start <= 1'b0;
            wr_start <= 1'b0;
            wr_data  <= 0;
        end else begin
            rd_start <= 1'b0;
            wr_start <= 1'b0;

            case (state)
                S_IDLE: begin
                    done <= 1'b0;
                    if (start_pulse) begin
                        n     <= 0;
                        state <= S_READ;
                    end
                end
                S_READ: begin
                    rd_start <= 1'b1;
                    state    <= S_WAIT_READ;
                end
                S_WAIT_READ: begin
                    if (a_done && b_done) begin
                        wr_data <= soma_sat;
                        state   <= S_WRITE;
                    end
                end
                S_WRITE: begin
                    wr_start <= 1'b1;
                    state    <= S_WAIT_WRITE;
                end
                S_WAIT_WRITE: begin
                    if (wr_done) begin
                        if (n == N_SAMPLES - 1) begin
                            state <= S_DONE;
                        end else begin
                            n     <= n + 1'b1;
                            state <= S_READ;
                        end
                    end
                end
                S_DONE: begin
                    done <= 1'b1;
                    if (!start) state <= S_IDLE;
                end
                default: state <= S_IDLE;
            endcase
        end
    end

    // --- controladores de memoria: os mesmos do conv1d e do iir_cascade ---
    memory_read_controller #(
        .DATA_WIDTH (DATA_WIDTH), .MEM_ADDRESS_N_BITS (ADDR_N_BITS)
    ) u_rd_a (
        .clk (clk), .reset_n (reset_n),
        .start_read (rd_start), .sram_readdata (a_sram_readdata),
        .data_addr (n), .mem_write (a_sram_write), .data_out (a_data),
        .mem_address (a_sram_address), .read_done (a_done),
        .clken (a_sram_clken), .chipselect (a_sram_chipselect)
    );

    memory_read_controller #(
        .DATA_WIDTH (DATA_WIDTH), .MEM_ADDRESS_N_BITS (ADDR_N_BITS)
    ) u_rd_b (
        .clk (clk), .reset_n (reset_n),
        .start_read (rd_start), .sram_readdata (b_sram_readdata),
        .data_addr (n), .mem_write (b_sram_write), .data_out (b_data),
        .mem_address (b_sram_address), .read_done (b_done),
        .clken (b_sram_clken), .chipselect (b_sram_chipselect)
    );

    memory_write_controller #(
        .DATA_WIDTH (DATA_WIDTH), .MEM_ADDRESS_N_BITS (ADDR_N_BITS)
    ) u_wr_y (
        .clk (clk), .reset_n (reset_n),
        .start_write (wr_start), .data_in_addr (n), .data_in (wr_data),
        .write_done (wr_done), .mem_address (y_sram_address),
        .mem_wdata (y_sram_wdata), .mem_write (y_sram_write),
        .clken (y_sram_clken), .chipselect (y_sram_chipselect)
    );

endmodule
\end{lstlisting}

% =============================================================================
\section{A \texttt{handle\_soma()} completa}
\label{ap:handle}
% =============================================================================

\begin{lstlisting}[language=C]
static int handle_soma(int sock, uint16_t dtype, uint32_t n_a, uint32_t n_b) {
    LOG("SOMA request: dtype=%u, n_a=%u, n_b=%u", dtype, n_a, n_b);

    if (!fpga_preparada()) return recusar_sem_fpga(sock, MORPHE_OP_SOMA, "SOMA");

    /* Sem esta conferencia, um servidor novo numa placa com o bitstream
     * antigo escreveria em 0x40000, onde nao ha memoria -- e acesso a
     * endereco sem escravo trava o barramento do HPS. */
    uint32_t ts = g_sysid[1];
    if (ts != SYSID_QSYS_TIMESTAMP) {
        char msg[192];
        snprintf(msg, sizeof msg,
                 "SOMA: o bitstream na FPGA (sysid %u) nao e o deste servidor (%u) -- "
                 "rode ./morphe-up.sh com o .sof que tem a soma", ts, SYSID_QSYS_TIMESTAMP);
        LOG("  -> %s", msg);
        return send_error(sock, MORPHE_OP_SOMA, MORPHE_STATUS_FPGA_NAO_PREPARADA, msg);
    }

    if (n_a == 0 || n_a > MORPHE_SOMA_N_MAX || n_b != n_a) {
        char msg[96];
        snprintf(msg, sizeof msg, "SOMA: a e b com o mesmo tamanho, de 1 a %d",
                 MORPHE_SOMA_N_MAX);
        return send_error(sock, MORPHE_OP_SOMA, MORPHE_STATUS_BAD_SIZE, msg);
    }
    uint32_t n = n_a;

    /* 1. recebe a e b */
    size_t payload_bytes = (size_t) 2 * n * 4;
    static uint8_t rx_buf[RX_BUF_MAX];
    if (payload_bytes > sizeof rx_buf)
        return send_error(sock, MORPHE_OP_SOMA, MORPHE_STATUS_BAD_SIZE, "SOMA: payload excedido");
    if (recv_exact(sock, rx_buf, payload_bytes) < 0) return -1;

    static int32_t a_buf[MORPHE_SOMA_N_MAX];
    static int32_t b_buf[MORPHE_SOMA_N_MAX];
    decode_samples_to_i32(rx_buf,         n, dtype, a_buf);
    decode_samples_to_i32(rx_buf + n * 4, n, dtype, b_buf);

    /* 2. escreve nas memorias da FPGA (porta s2), completando com zeros */
    for (uint32_t i = 0; i < MORPHE_SOMA_N_MAX; i++) g_soma_a[i] = (i < n) ? a_buf[i] : 0;
    for (uint32_t i = 0; i < MORPHE_SOMA_N_MAX; i++) g_soma_b[i] = (i < n) ? b_buf[i] : 0;

    /* 3. dispara (borda de subida no start) e espera o done */
    *g_pio_soma_start = 0;
    usleep(1);
    *g_pio_soma_start = 1;
    if (wait_done(g_pio_soma_done, FPGA_DONE_TIMEOUT_MS) < 0) {
        *g_pio_soma_start = 0;
        return send_error(sock, MORPHE_OP_SOMA, MORPHE_STATUS_FPGA_TIMEOUT, "SOMA: timeout");
    }
    *g_pio_soma_start = 0;      /* com start em 0, o soma.v volta ao repouso */

    /* 4. le a saida */
    static int32_t y_buf[MORPHE_SOMA_N_MAX];
    for (uint32_t i = 0; i < n; i++) y_buf[i] = g_soma_y[i];

    /* o mesmo .mrph de depuracao da CONV: a no lugar de x, b no de h */
    save_debug_bundle_conv("soma", dtype, n, a_buf, n, b_buf, n, y_buf);

    /* 5. responde */
    static uint8_t tx_buf[TX_BUF_MAX];
    uint8_t hdr[MORPHE_HEADER_SIZE];
    build_resp_header(hdr, MORPHE_OP_SOMA, dtype, MORPHE_STATUS_OK, n);
    if (dtype == MORPHE_DTYPE_INT32) {
        for (uint32_t i = 0; i < n; i++) i32_to_be(tx_buf + i * 4, y_buf[i]);
    } else {
        for (uint32_t i = 0; i < n; i++) f32_to_be(tx_buf + i * 4, (float) y_buf[i]);
    }
    if (send_all(sock, hdr, sizeof hdr) < 0) return -1;
    if (send_all(sock, tx_buf, (size_t) n * 4) < 0) return -1;

    LOG("  -> SOMA OK: n=%u", n);
    return 0;
}
\end{lstlisting}

% =============================================================================
\section{A janela \texttt{soma\_window.py} completa}
\label{ap:janela}
% =============================================================================

\begin{lstlisting}[language=Python]
import threading
import tkinter as tk
from tkinter import ttk, messagebox
from typing import Optional

import numpy as np

import matplotlib
matplotlib.use("TkAgg")
from matplotlib.figure import Figure
from matplotlib.backends.backend_tkagg import FigureCanvasTkAgg

import dsp_core as dsp
import morphe_theme as theme
from conv_window import _stem_signal
from morphe_config import SOMA_N_MAX
from morphe_protocol import DTYPE_INT32, build_soma_request, decode_soma_response
from plot_toolbar import PlotToolbar
from signal_panel import SignalPanel


def soma_referencia(a_q: np.ndarray, b_q: np.ndarray) -> np.ndarray:
    """A conta do soma.v em NumPy: soma em 64 bits e satura na faixa de 32."""
    s = np.asarray(a_q, dtype=np.int64) + np.asarray(b_q, dtype=np.int64)
    return np.clip(s, dsp.Q_MIN, dsp.Q_MAX)


class SomaWindow(tk.Toplevel):
    def __init__(self, master):
        super().__init__(master)
        self.title("Morphe — Soma via FPGA (exemplo do roteiro)")
        self.geometry("1240x820")
        self.configure(bg=theme.COLORS["bg"])
        theme.setup_styles(self)

        self.a_sig: Optional[dsp.Signal] = None
        self.b_sig: Optional[dsp.Signal] = None
        self.y_sig: Optional[dsp.Signal] = None

        self.status = tk.StringVar(value="Gere a[n] e b[n] e clique em Somar na FPGA.")
        theme.make_status_bar(self, self.status).pack(side="bottom", fill="x")

        root = ttk.Frame(self, style="Main.TFrame", padding=12)
        root.pack(fill="both", expand=True)
        lado = ttk.Frame(root, style="Main.TFrame")
        lado.pack(side="left", fill="y", padx=(0, 12))
        self._build_sidebar(lado)
        area = ttk.Frame(root, style="Main.TFrame")
        area.pack(side="right", fill="both", expand=True)
        self._build_plots(area)
        self._redraw()

    def _build_sidebar(self, parent):
        self.a_panel = SignalPanel(parent, title="Sinal a[n]", on_generate=self._on_gen_a,
                                   default_N=32, default_type="Senóide", max_N=SOMA_N_MAX)
        self.a_panel.pack(fill="x")
        self.b_panel = SignalPanel(parent, title="Sinal b[n]", on_generate=self._on_gen_b,
                                   default_N=32, default_type="Degrau unitário", max_N=SOMA_N_MAX)
        self.b_panel.pack(fill="x", pady=(8, 0))

        op = ttk.LabelFrame(parent, text=" Operação ", style="Card.TLabelframe",
                            padding=(14, 12, 14, 14))
        op.pack(fill="x", pady=(12, 0))
        ttk.Label(op, style="SectionHint.TLabel", wraplength=300, justify="left",
                  text=f"Até {SOMA_N_MAX} amostras, em Q15.16. Se os tamanhos forem "
                       "diferentes, o menor é completado com zeros.").pack(anchor="w", pady=(0, 8))
        self.btn = ttk.Button(op, text="Somar na FPGA", style="Primary.TButton",
                              command=self._on_somar)
        self.btn.pack(fill="x")

    def _build_plots(self, parent):
        self.fig = Figure(figsize=(7, 8), dpi=100)
        self.fig.patch.set_facecolor(theme.COLORS["bg"])
        self.ax_a = self.fig.add_subplot(311)
        self.ax_b = self.fig.add_subplot(312)
        self.ax_y = self.fig.add_subplot(313)
        self.canvas = FigureCanvasTkAgg(self.fig, master=parent)
        self.canvas.get_tk_widget().pack(fill="both", expand=True)
        self.toolbar = PlotToolbar(self.canvas)

    def _redraw(self):
        _stem_signal(self.ax_a, self.a_sig, "a[n]", dsp.COLOR_X)
        _stem_signal(self.ax_b, self.b_sig, "b[n]", dsp.COLOR_H)
        _stem_signal(self.ax_y, self.y_sig, "y[n] = a[n] + b[n] (FPGA)", dsp.COLOR_Y)
        self.fig.tight_layout()
        self.canvas.draw_idle()

    def _on_gen_a(self, panel: SignalPanel):
        try:
            self.a_sig = panel.build()
            self._redraw()
        except Exception as e:
            messagebox.showerror("Erro em a[n]", str(e))

    def _on_gen_b(self, panel: SignalPanel):
        try:
            self.b_sig = panel.build()
            self._redraw()
        except Exception as e:
            messagebox.showerror("Erro em b[n]", str(e))

    def _on_somar(self):
        if self.a_sig is None or self.b_sig is None:
            messagebox.showwarning("Atenção", "Gere a[n] e b[n] antes.")
            return
        try:
            client = self.master.tcp_panel.make_client()
        except Exception as e:
            messagebox.showerror("Configuração", str(e))
            return

        # mesmo tamanho para os dois; o indice n vem de a[n]
        n = max(self.a_sig.x.size, self.b_sig.x.size)
        a_q = dsp.float_to_q1516(dsp.pad_zeros_to(self.a_sig.x, n))
        b_q = dsp.float_to_q1516(dsp.pad_zeros_to(self.b_sig.x, n))
        n0 = int(self.a_sig.n[0]) if self.a_sig.n.size else 0

        self.btn.config(state="disabled")
        self.status.set(f"Enviando {n} amostras de a e de b à FPGA...")

        def worker():
            try:
                resp = client.request(build_soma_request(a_q, b_q, DTYPE_INT32))
                y_q = decode_soma_response(resp)
                self.after(0, lambda: self._on_resultado(y_q, a_q, b_q, n0))
            except Exception as e:
                self.after(0, lambda err=e: self._on_erro(err))

        threading.Thread(target=worker, daemon=True).start()

    def _on_resultado(self, y_q, a_q, b_q, n0):
        ref = soma_referencia(a_q, b_q)
        iguais = int(np.sum(y_q == ref))
        n = y_q.size
        self.y_sig = dsp.Signal(n=np.arange(n) + n0, x=dsp.q1516_to_float(y_q),
                                description=f"{n} amostras, Q15.16")
        self._redraw()
        self.btn.config(state="normal")
        if iguais == n:
            self.status.set(f"OK: as {n} amostras da FPGA são iguais, bit a bit, "
                            "à soma em ponto fixo feita no NumPy.")
        else:
            self.status.set(f"ATENÇÃO: só {iguais} de {n} amostras iguais à referência.")

    def _on_erro(self, err: Exception):
        self.btn.config(state="normal")
        self.status.set("Erro — ver mensagem.")
        messagebox.showerror("Erro na soma", str(err))
\end{lstlisting}

\end{document}
